\documentclass[12pt]{article}
\usepackage[utf8]{inputenc}
\usepackage[margin=0.75in]{geometry}
\usepackage{amsmath}\usepackage{amssymb}\usepackage{enumitem}\usepackage{multicol}\usepackage{tcolorbox}
\setlength{\parindent}{0pt}
\begin{document}

\begin{center}{\Large \textbf{State the domain and range of a function}}\\[2pt]{\small Pre-Cal \textbar\ CK-12: Domain and Range}\end{center}\vspace{0.25cm}

{\large\textbf{Key Idea}}\par\smallskip
\textbf{Domain} is every allowed input $x$; \textbf{range} is every output $y$. Two things restrict a domain: you cannot divide by 0, and you cannot take the square root of a negative.\par


{\large\textbf{Key Terms}}\par\smallskip
\begin{itemize}[leftmargin=1.4em]
  \item \textbf{Domain} — allowed $x$-values
  \item \textbf{Range} — resulting $y$-values
\end{itemize}


{\large\textbf{Worked steps}}\par\smallskip
Step 1: Set any denominator $\neq 0$; set any radicand $\ge 0$.\par
Step 2: Solve for $x$ to get the domain.\par
Example: for $f(x)=\sqrt{x-2}$, need $x-2\ge 0$, so $x\ge 2$.\par


{\large\textbf{You Try}}\par\smallskip
Give the domain of $f(x)=\dfrac{5}{x+4}$.\par

\end{document}